\section*{Chapter \ref{ch:b3:planetary-interiors-surfaces} --- Planetary Interiors and Surfaces}

\begin{solution}{exo:b3:planetary-interiors-surfaces:1}
$\tfrac52C\bar\rho = 2.5\times0.35\times5429 = 4750$; $x^2 = (4750 -
3300)/(5429 - 3300) = 0.681$, $x = 0.825$, a core of $0.825\times2440 =
\qty{2010}{km}$; $\rho_c = 3300 + 2129/0.825^3 = \qty{7090}{kg.m^{-3}}$.
\end{solution}

\begin{solution}{exo:b3:planetary-interiors-surfaces:2}
$m = 3500\times\tfrac43\pi(25)^3 = \qty{2.29e8}{kg}$; $E = \tfrac12 m v^2 =
\tfrac12\times\num{2.29e8}\times(\num{2.0e4})^2 = \qty{4.6e16}{J}$, or
$\num{4.6e16}/\num{4.2e15} = \qty{11}{Mt}$. The rule gives a crater of
$20\times\qty{50}{m} = \qty{1}{km}$.
\end{solution}

\begin{solution}{exo:b3:planetary-interiors-surfaces:3}
$N_1 = 120/\num{100000} = \qty{1.2e-3}{km^{-2}}$, relative error $1/\sqrt{120}
= \qty{9}{\percent}$. On the linear part, $T = N_1/c = \num{1.2e-3}/\num{4.875e-4}
= \qty{2.46}{Ga}$ (the exponential term, $\num{8e-7}$, is negligible), and
$\sigma_T = T/\sqrt{120} = \qty{0.22}{Ga}$.
\end{solution}

\begin{solution}{exo:b3:planetary-interiors-surfaces:4}
Area $4\pi(\num{1.8215e6})^2 = \qty{4.17e13}{m^2}$; power $2.5\times
\num{4.17e13} = \qty{1.04e14}{W} = \qty{104}{TW}$, 2.2 times the Earth's
total from a surface twelve times smaller, and $2.5/0.092 = 27$ times its
flux.
\end{solution}

\begin{solution}{exo:b3:planetary-interiors-surfaces:5}
(a)~$\bar\rho = \num{6.417e23}/(\tfrac43\pi(\num{3.3895e6})^3) =
\qty{3934}{kg.m^{-3}}$; $g = GM/R^2 = \qty{3.73}{m.s^{-2}}$; $\rho gR =
\qty{4.97e10}{Pa}$. (b)~$19\mu/(2\rho gR) = 19\times\num{6.7e10}/\num{9.94e10}
= 12.8$, $k_2 = 1.5/13.8 = 0.109$. (c)~The measured value is larger by
2.6 standard deviations: Mars yields more than a solid planet with a
mantle's rigidity, as it would with a liquid core inside.
\end{solution}

\begin{solution}{exo:b3:planetary-interiors-surfaces:6}
(a)~Moon: $4\times9.8/1.62 = \qty{24}{km}$ in crystalline rock,
\qty{12}{km} in sediments; Mars and Mercury: $4\times9.8/3.7 =
\qty{10.6}{km}$ and \qty{5.3}{km}. The \qty{20}{km} of the Moon lies between
its two predictions: the rule holds to a few tens of percent. (b)~$Y =
\rho gD/3 = 2600\times9.8\times4000/3 = \qty{34}{MPa}$, the strength of
fractured rock rather than of an intact crystal.
\end{solution}

\begin{solution}{exo:b3:planetary-interiors-surfaces:7}
Mars: $\num{3.22e-14}(e^{6.93T} - 1) + \num{4.875e-4}\,T = \num{2.0e-3}$;
the linear guess $4.1$~Ga is far too old, the exponential being large
there; bisection gives $T = \qty{3.34}{Ga}$ (exponential term
$\num{3.6e-4}$, linear term $\num{1.63e-3}$). The lunar function gives
\qty{2.39}{Ga}. Each function is calibrated for one body: the rate of
impacts and the size of the craters a given projectile makes differ
between the Moon and Mars.
\end{solution}

\begin{solution}{exo:b3:planetary-interiors-surfaces:8}
(a)~$3.551181/1.769138 = 2.0073$ and $7.154553/3.551181 = 2.0147$, close to
$2$. (b)~$0.565247 - 0.844789 + 0.279542 = 0.000000$ per day, to
$\num{1e-7}$: the three mean motions are locked, a resonance that holds
the moons' relative positions at each conjunction. (c)~The pulls repeat at
the same points of the orbit, so they add up instead of averaging out;
they feed the eccentricity as fast as the tides damp it.
\end{solution}

\begin{solution}{exo:b3:planetary-interiors-surfaces:9}
$n = 2\pi/(1.769138\times86400) = \qty{4.111e-5}{s^{-1}}$, $nR =
\qty{74.9}{m.s^{-1}}$, $(nR)^5/G = \qty{3.53e19}{W}$;
$\tfrac{21}{2}e^2(nR)^5/G = 10.5\times\num{1.6e-5}\times\num{3.53e19} =
\qty{5.9e15}{W}$ per unit of $k_2/Q$. Then $k_2/Q =
\num{1.04e14}/\num{5.9e15} = 0.018$.
\end{solution}

\begin{solution}{exo:b3:planetary-interiors-surfaces:10}
(a)~$x^2 = (\tfrac52C\bar\rho - \rho_m)/(\bar\rho - \rho_m)$ is linear in
$C$, of slope $\tfrac52\bar\rho/(\bar\rho - \rho_m)$. (b)~The slope is
$2.5\times3340/52 = 161$: $x^2$ changes by $0.161$ for $0.001$ of $C$.
For $C = 0.395$, $x^2 = 0.036 + 0.161 = 0.197$, $x = 0.44$, $\rho_c =
\qty{3880}{kg.m^{-3}}$; for $C = 0.393$, $x^2 < 0$: no core at all.
(c)~Near $C = \tfrac25$ the moment of inertia barely feels a small core:
its third significant figure decides between a large light core and none.
Seismic waves, tides and the magnetic field are needed.
\end{solution}

\begin{solution}{exo:b3:planetary-interiors-surfaces:11}
(a)~$\sigma_T = N_1/(10\,N_1')$: at \qty{1.0}{Ga}, $N_1 \approx N_1' \approx
\num{4.875e-4}$, $\sigma_T = \qty{0.10}{Ga}$; at \qty{3.8}{Ga}, $N_1 =
\num{1.07e-2}$ and $N_1' = \num{6.15e-2}$, $\sigma_T = \qty{0.017}{Ga}$,
six times smaller. (b)~With $N_1 = cT$, $N_1' = c$ and $\sigma_T =
cT/(\sqrt n\,c) = T/\sqrt n$. (c)~The chronology function itself is
uncertain: it rests on a few dated sites of the Moon and on the scaling
from the Moon to Mars, errors that no count reduces.
\end{solution}

\begin{solution}{exo:b3:planetary-interiors-surfaces:12}
(a)~The ratio is $(n_ER_E/n_IR_I)^5(e_E/e_I)^2 = (0.4982\times0.8569)^5
\times(0.009/0.004)^2 = 0.0142\times5.06 = 0.072$: about
\qty{7.5e12}{W}, or $\num{7.5e12}/(4\pi(\num{1.5608e6})^2) =
\qty{0.24}{W.m^{-2}}$. (b)~An ocean decouples the ice shell from the
interior: the shell flexes far more than a solid moon would, and its ice
dissipates differently from Io's hot rock, so $k_2$ and $Q$ both change.
\end{solution}

\begin{solution}{pb:b3:planetary-interiors-surfaces:1}
\textbf{1.} $\bar\rho = \num{6.4169e23}/(\tfrac43\pi(\num{3.3895e6})^3) =
\qty{3934}{kg.m^{-3}}$.

\textbf{2.} $\tfrac52C\bar\rho = 3600$; $x^2 = (3600 - 3500)/(3934 - 3500) =
0.230$, $x = 0.479$, a core of \qty{1620}{km}; $\rho_c = 3500 +
434/0.479^3 = \qty{7450}{kg.m^{-3}}$.

\textbf{3.} $x = 1830/3389.5 = 0.540$; $\rho_m = (3600 - 0.2915\times
3934)/(1 - 0.2915) = \qty{3460}{kg.m^{-3}}$, $\rho_c = 3462 +
472/0.1574 = \qty{6460}{kg.m^{-3}}$.

\textbf{4.} It is a third lower. The pressure in the core of Mars is far
lower than in the Earth's, so its metal is less compressed; and a core so
light suggests much lighter elements mixed with the iron.

\textbf{5.} The \omterm{def:b3:planetary-interiors-surfaces:love-number}{Love number} (2003) came long before InSight (2021): it
already showed that at least the outer core is liquid and placed its
radius between 1520 and \qty{1840}{km}, a range that contains the
seismological value.

\textbf{6.} $1 - (1 - 20/3389.5)^3 = \qty{1.8}{\percent}$ to $1 - (1 -
37/3389.5)^3 = \qty{3.2}{\percent}$ of the volume, and less of the mass:
lighter than the mantle, it lowers the effective mantle density a little,
an effect smaller than the uncertainty of the model.

\textbf{7.} From \qty{1}{km}: 54; from 1.41: 16; from 2: 10; from 2.83: 5;
from 4, 5.66, 8 and 11.3: 2; from 16: 1.

\textbf{8.} $N = 54/80000 = \qty{6.75e-4}{km^{-2}}$, $\sigma = \sqrt{54}/80000
= \qty{0.92e-4}{km^{-2}}$ (\qty{14}{\percent}).

\textbf{9.} Expected $54\times0.355 = 19$ craters from \qty{1.41}{km} and
$54\times0.147 = 7.9$ from \qty{2}{km}; counted 16 and 10, within one
standard deviation ($\sqrt{19} = 4.4$, $\sqrt{7.9} = 2.8$): consistent.

\textbf{10.} No: $0.072/\num{6.75e-4} = 107$ times below it. The count
records time.

\textbf{11.} Bins of equal width in $\log D$ cover the range evenly and
keep comparable numbers of craters in the bins of a steep distribution.

\textbf{12.} They are secondary craters, made by blocks thrown out of a
large primary crater: not independent impacts from space, they would add
craters the chronology function does not count.

\textbf{13.} $T = \num{6.75e-4}/\num{4.875e-4} = \qty{1.385}{Ga}$; the
exponential term, $\num{3.22e-14}(e^{9.6} - 1) = \num{4.7e-10}$, is a
millionth of $N_1$.

\textbf{14.} $\sigma_T = T/\sqrt{54} = \qty{0.19}{Ga}$: $1.38 \pm 0.19$~Ga.

\textbf{15.} $\num{6.75e-4}/\num{8.38e-4} = \qty{0.81}{Ga}$, younger by
almost half.

\textbf{16.} $31 + 12 + 7 + 2 + 0 + 2 + 2 = 56$ craters; $N(\qty{2}{km}) =
56/120000 = \qty{4.67e-4}{km^{-2}}$; $N_1 = \num{4.67e-4}/0.1473 =
\qty{3.17e-3}{km^{-2}}$. Bisection on $\num{3.22e-14}(e^{6.93T} - 1) +
\num{4.875e-4}\,T = \num{3.17e-3}$ gives $T = \qty{3.54}{Ga}$ (exponential
term \num{1.45e-3}, linear term \num{1.73e-3}).

\textbf{17.} $N_1' = \num{3.22e-14}\times6.93\,e^{6.93\times3.54} +
\num{4.875e-4} = \num{1.05e-2}$; $\sigma_T = \num{3.17e-3}/(\sqrt{56}\times
\num{1.05e-2}) = \qty{0.040}{Ga}$.

\textbf{18.} Plain A: \qty{14}{\percent}; plain B: \qty{1.1}{\percent}. In
the steep, early part of the chronology a large change of the crater
density means a small change of age; on the linear part the age is as
uncertain as the count.

\textbf{19.} $120000\times\num{3.17e-3}\times(1 - 0.1473) = 324$; it
carries $48 + 8 = 56$.

\textbf{20.} $N_1 = 56/(120000\times0.8527) = \qty{5.47e-4}{km^{-2}}$,
$T_r = \num{5.47e-4}/\num{4.875e-4} = \qty{1.12}{Ga}$.

\textbf{21.} $\sigma = 1.12/\sqrt{56} = \qty{0.15}{Ga}$.

\textbf{22.} A \qty{2}{km} crater is \qty{400}{m} deep: the lava was at
least that thick. Larger craters are deeper and their rims stand higher;
they were partly filled and still show.

\textbf{23.} $1.38 \pm 0.19$ and $1.12 \pm 0.15$~Ga differ by \qty{0.26}{Ga},
about one combined standard deviation ($\sqrt{0.19^2 + 0.15^2} = 0.24$):
the counts allow a single episode, though they do not prove it.

\textbf{24.} The young lava plain A has a model age of
\textbf{1.4~billion~years} ($1.38 \pm 0.19$~Ga).
\end{solution}
